\chapter{Discussion}\label{chapter:discussion}

This chapter interprets the results of the thesis and places them within the
scope established by the research questions. It first discusses what the
feature model, reference architecture, and evaluation show when considered
together. It then examines the boundaries of the prototype and the evaluation,
followed by directions for future work.

\section{Interpretation of the Results}

\subsection{From the Monitoring Design Space to a Reusable Infrastructure}

The review in Chapter~\ref{chapter:related_work} shows that existing approaches
usually specialize in one part of the monitoring problem. Some approaches
collect ROS messages for an external oracle, some generate checkers from
requirements, and others observe network traffic or lower-level execution
events. Topic monitoring is common in the reviewed ROS-focused approaches.
ROSMonitoring 2.0 extends this coverage to services, with partial ROS~2
support, while ROS~2 action monitoring remains an open gap in the reviewed
work. No reviewed approach combines typed ROS~2 topic messages, service events,
and action progress in a configurable infrastructure that can place collection
and checking in separate processes. These boundaries explain the observation
gaps identified for RQ1 and the broader research gap described in
Chapter~\ref{chapter:related_work}.

The feature model developed for RQ2 gives structure to this broad design space.
Its main value is the separation of observable elements, monitoring
architecture, monitoring mode, system configuration, and implementation
paradigm. These dimensions make the design choices explicit. For example, the
choice between application-level and lower-level observation is separate from
the choice between integrated and distributed placement. This separation
prevents one deployment decision from becoming a fixed property of every
collector or checker.

The prototype implements the focused path summarized in
Table~\ref{tab:selected-feature-coverage}: passive application-level collection
through ROS~2-facing monitor nodes, configurable processing and checking, and
integrated or split deployment. The selected path covers the research gap
identified in Chapter~\ref{chapter:related_work}. Topic messages, service
events, and action feedback and status enter the same pipeline, while
converters and verdict services keep rule-specific knowledge outside the
collection layer. This result answers RQ3 for the selected feature set. The reference
architecture incorporates the common monitoring tasks and preserves the
variation points needed for sources, processing steps, checking components,
outputs, and placement.

Several unselected feature variants map onto these variation points and could
therefore be hosted without structural change. Parameter changes and
lifecycle transitions are published on regular topics
(\texttt{/parameter\_events} and the lifecycle transition-event topics), so
they need only additional collector configurations or collector types. A
complete action trace requires service collectors for the goal, cancellation,
and result events of the hidden action services. A new transport requires an
exporter and a matching source, and a new checker a converter and
verdict-service pair. Other variants would change the assumptions of the
design: active intervention needs a feedback path into the application,
global message ordering and clock guarantees need coordination that the
record format alone does not provide, and network taps or kernel tracing
observe data without the typed source identity that the collectors rely on.

The common data record is central to this result. Uniform fields let sources,
transformers, transports, and checking components exchange one representation.
At the same time, \texttt{source\_type}, \texttt{source\_name},
\texttt{phase}, timestamps, and metadata retain the distinctions between ROS~2
communication mechanisms. The uniform representation therefore simplifies
the pipeline boundary while preserving the context needed by later checking.
The session and record identifiers also provide an evidence trail from a
verdict back to its contributing observations. E1--E5 use this attribution to
compare online verdicts with recorded inputs.

This architecture reduces repeated monitoring work. Runtime type resolution,
message conversion, record creation, JSONL serialization, routing, and
transport are implemented once in the infrastructure. A new application still
requires decisions about useful sources and rules. Once these decisions
are made, YAML configuration and project-specific plugins reuse the same
monitoring core. The prototype therefore automates the recurring
infrastructure tasks described in the problem statement while leaving
application knowledge explicit.

\subsection{Meaning of the Evaluation Results}

The five experiments provide complementary evidence for RQ4. E1 establishes
the complete integrated path for a topic rule. E2 shows that topic
messages, service request and response events, and action feedback and status
can coexist in one session. It also shows that a rule can correlate a
service response with later topic data. Its identical local and remote speed
verdicts show that serialization and MQTT transport preserved the meaning and
record attribution of the check in that run.

E3 is important because it replaces the small stimulus processes with an
existing Nav2 application. The monitor observes Nav2 through its published
interfaces, and the navigation code remains unchanged. The three goals cover a
\textsc{succeeded} status within the deadline, a goal that remains active beyond
the deadline, and an \textsc{aborted} status. Agreement between the online
verdicts, the mission log, and the Gazebo route data shows that the checking
chain responded to meaningful changes in application behaviour. It also
demonstrates the value of combining topic messages with action feedback and
status: topic data describes robot motion, while the action records describe
progress and the final status.

E4 extends this result from one trace to two robot namespaces. The independent
speed checks show that \texttt{source\_name} keeps the robot streams separate.
The separation check shows that one converter can combine observations from
both robots. This is a useful form of multi-robot monitoring because several
robot-local streams can feed one rule while the collection mechanism
remains unchanged. The agreement with Gazebo positions gives an additional
source of evidence for the four separation transitions.

E5 changes placement while retaining the E4 checking graph. DDS carries the
robot observations from the workstation to the Raspberry Pi, and MQTT carries
the data records to a ROS-independent verifier on macOS. Both verifier
processes received every data record from the six monitored E4 sources, and
the 34 verdict transitions on the macOS verifier matched the Pi reference.
These results show that the split deployment works across the operating systems
and processor architectures used in the experiment. They also support the
design decision to keep ROS~2 dependencies in the collection layer: the remote
verifier needs the record and checking interfaces, while ROS~2 type handling
stays near the system under monitoring.

Taken together, the experiments support a positive answer to RQ4 within the
scope stated in Section~\ref{sec:eval-rq4-answer}. The infrastructure handles
the selected application-level observations, processing steps, checking
placements, and multi-source rules in both small examples and Nav2 case
studies. The results establish feasibility and functional correctness for the
observed runs. Broader guarantees depend on the limitations discussed below.

\subsection{Design Trade-offs}

Early processing is useful for high-rate ROS~2 sources. Field extraction keeps
only the values required by a rule, rate limiting reduces repeated
observations, and change filtering suppresses stable values. E3 demonstrates
this choice by reducing odometry before storage and MQTT transport. The same
choice makes the monitoring configuration part of the checking setup. A
rate limit can hide a short-lived state, and field extraction can remove data
needed by a later checker. Users therefore need to select transformers with the
timing and data needs of each rule in mind. Keeping the applied transformer
names in record metadata helps later interpretation. Discarded observations
remain unavailable to downstream components.

The split deployment moves checking work away from the ROS~2 application and
allows the verifier to use a plain Python environment. It also introduces a
transport boundary. Network delay, temporary disconnection, queue limits, and
broker availability can affect when records reach the checker. The integrated
deployment avoids this boundary and keeps the trace within one process, while
sharing processor and memory resources with the ROS~2-facing monitor. The two
deployment forms therefore serve different operational needs. The
configuration keeps this choice separate from converter and verdict-service
logic, which allows the same checking graph to be placed according to the
available resources and network.

Plugin interfaces provide flexibility at the checking boundary. A user can
provide a converter and a verdict service as Python classes and select them in
YAML by their import paths. Together, these two plugins form an adapter for the
user's DSL. The converter changes the common data record into the input format
of that DSL. The verdict service holds or calls the checking engine and changes
its result into the common verdict format. The infrastructure treats the data
between these plugins as a DSL record, so the DSL syntax and checking logic stay
inside the user-provided classes.

The converter and verdict service can run in the same process or on different
hosts. In the second case, the DSL record is serialized as JSON and sent through
a file or MQTT link. This keeps the DSL adapter separate from its placement,
provided that the converter produces a JSON-serializable record.

The case-study plugins show several simple uses of this extension point. A
converter can select one value, combine records from several sources, or emit a
record when a timer expires. A verdict service can check a threshold, a
deadline, or a change of state. The same interfaces can also connect a larger
DSL engine. For example, a converter could prepare timestamped signals for a
temporal-logic DSL, or translate ROS~2 observations into events for a
project-specific state-machine DSL. These extensions reuse the common record,
checking graph, deployment configuration, and verdict exporters.

\section{Limitations of the Prototype}

\subsection{Observation and Checking Scope}

The prototype concentrates on passive application-level observation. A topic
collector sees the messages delivered to its own subscription, so compatible
QoS settings and successful DDS discovery remain necessary. Service monitoring
depends on service introspection being enabled by a client or server. Action
collection covers feedback and status directly. Goal, cancellation, and result
exchanges require observation of the corresponding service events, which means
that a complete action trace needs service collectors in addition to the
current action collector.

Several feature-model variants remain outside this implementation. These
include parameters, node lifecycle states, DDS-level QoS events, network
packets, and operating-system or tracing events. The prototype also uses
ROS~2-facing nodes as its observation mechanism. Decorator-based interception,
network tapping, and LTTng tracing would require additional collector types and
different assumptions about types and source identity.

The implemented monitoring mode is passive. Verdicts are exported to standard
output, files, or MQTT for later use. Recovery actions, message blocking, and
changes to robot behaviour remain outside the runtime path. This scope keeps
the monitor from directly affecting application decisions, while limiting its
role to observation, checking, and reporting.

The infrastructure leaves the choice of rule language to user-provided
plugins. Users can support their own DSL by implementing a converter and a
verdict service, and the framework checks that both classes implement the
required interfaces. The two plugins still have to agree on the shape of the
DSL record, and the user has to implement and test the meaning of the rule.
The current project provides examples of this process, while reusable
templates, common test support, and clearer plugin descriptions could make
custom DSL integration easier.

Global rules are currently implemented by converters that keep the
required state from several sources. Automatic decomposition of a rule
among several monitors is beyond the current split deployment, whose checking
remains logically centralized at a verifier.

\subsection{Ordering, Time, and Delivery}

Each data record has a monitor-side wall-clock timestamp and a per-source
sequence number. These fields support attribution and ordering within one
collector. An integrated monitor also has a local processing order for the
records it receives. The infrastructure does not reconstruct publication order
across several publishers or establish a global order across hosts and clocks.
A distributed rule can therefore observe records in an order influenced by
DDS and MQTT delivery. The E4 and E5 separation converter pairs each new
position with the latest position received for the other robot. This approach
is suitable for the evaluated rule, while its result contains the sampling
and arrival-time difference between the two streams.

MQTT QoS and record identifiers provide useful transport support. End-to-end
exactly-once delivery remains outside the current prototype. A queued message
can be dropped when a bounded outbound queue is full, and QoS~1 delivery can
introduce duplicates after reconnection. Record identifiers make such
duplicates recognizable. Duplicate suppression, persistent buffering,
recovery of missed intervals, and a defined policy for late records require
further work.

The prototype also lacks fault-tolerant checking and automatic failover. A
broker or verifier failure can interrupt a split monitoring run, and a
logically centralized verifier can become a bottleneck as the number or rate of
sources grows. Authentication, access control, and protection of monitoring
data were outside the evaluation. These concerns become more important when a
split deployment crosses an untrusted network or carries sensitive robot data.

\subsection{Configuration and Operation}

The deployment generator reduces inconsistency between host configurations by
projecting one JSON deployment description into matching runtime YAML files.
The description still requires explicit hosts, runtime roles, sources, class
paths, links, and transport parameters. Users need knowledge of the ROS graph,
message fields, checker inputs, and deployment network before they can create a
valid setup. Configuration validation catches structural errors, while the
meaning of a selected source or rule still depends on the application.

Source types can be discovered as a startup fallback, and the required ROS~2
interface packages must be available in the monitor environment. The running
pipeline uses the predefined collectors and component placement for the full
run. Automatic reconfiguration after ROS graph or host changes remains future
work. Systems with dynamic nodes and long-lived operation also need clearer
operational state reporting.

\section{Threats to Validity}

The answers to RQ1 and RQ2 depend on the selection and interpretation of the
related work. The reviewed approaches cover several observation levels and
checking styles, while additional research or industrial systems may expose
further features. Differences in terminology and the level of detail provided
by each publication can also affect the capability classification. The
prototype shows that the selected path through the feature model can be
implemented. Independent review of the model by ROS~2 developers and monitoring
experts remains future validation work.

The evaluation combines deterministic stimuli, recorded JSONL artifacts,
offline replay, mission logs, and Gazebo positions. This combination strengthens
the correctness argument because verdicts can be checked against both recorded
inputs and application-level outcomes. Some comparisons replay the same records
through rule logic that is closely related to the online checker. A shared
implementation error could therefore affect both results. The independently
recorded mission and Gazebo data in E3--E5 reduce this threat for navigation
outcomes and robot separation. Independent oracle coverage remains limited for
individual field conversions and pipeline operations.

The experiments select representative configurations from a larger feature
model. They cover three ROS~2 communication mechanisms, two deployment forms,
multi-source checks, and a three-machine setup. Many valid combinations remain
untested, including other plugins, file transport across hosts, different QoS
profiles, and several simultaneous verifier chains. The results therefore
provide evidence for the implemented paths exercised by E1--E5. Other
configurations admitted by the model require separate evaluation.

External validity is limited by the experimental environment. The case studies
use ROS~2 Kilted, Fast DDS, one MQTT broker, one local network, at most two
simulated TurtleBot3 robots, and runs lasting a few minutes. E5 distributes the
runtime over three physical computers, while Gazebo remains the source of the
robot behaviour. Other ROS~2 distributions, DDS implementations, robot
hardware, networks, fleet sizes, and long-running deployments may expose
different timing, discovery, and resource behaviour.

The evaluation focuses on whether the infrastructure supports the selected
monitoring needs and whether the observed verdicts are correct. A systematic
scalability study, failure-injection study, usability study, and performance
comparison with the approaches in Chapter~\ref{chapter:related_work} remain
future evaluation work. The five experiments provide limited evidence about
production capacity and relative performance.

\section{Future Work}

The first direction is broader ROS~2 observation. A complete action trace could
combine feedback and status with goal, cancellation, and result service events.
Additional collectors could cover parameters, lifecycle transitions, and
DDS-level QoS events. Adapters for \texttt{ros2\_tracing} or network-level data
could complement the typed application records when a rule needs callback,
scheduling, middleware, or security information. The common data record would
need clear extensions for the identity and timing information of these sources.

The second direction is stronger support for distributed monitoring. Record
handling could add duplicate suppression, persistent queues, replay after
reconnection, and explicit rules for records that arrive late. Clock
synchronization and time windows based on record timestamps could improve
rules that correlate observations from several hosts. Larger deployments
would also benefit from several verifier
processes, automatic placement, health checking, and recovery after component
failure. Long-duration experiments over less reliable networks and with larger
robot fleets would show how these mechanisms affect delivery, latency, and
resource use. Security should be evaluated together with these extensions.

The converter and verdict-service interfaces already allow users to add their
own DSL and checking engine. Future work can make this extension path easier to
use by providing plugin templates, a small test harness, and clear documentation
of plugin parameters and DSL-record fields. Additional plugin packages could
then support temporal-logic DSLs, state-machine DSLs, or project-specific rule
languages. In each case, the converter would prepare the DSL input and keep the
input record identifiers, while the verdict service would run the checker and
return the common verdict. The collection and deployment parts of the
infrastructure would remain unchanged.

Exporters and transport endpoints provide a similar extension point. The
interfaces already accept user-defined exporter and source classes. The
built-in transports for data records and DSL records currently use files or
MQTT, and verdicts can also be written to standard output. Future work could
add ready-to-use plugins for ROS~2 topics, TCP or other socket connections,
HTTP or WebSocket endpoints, databases, or other message brokers. A new output
destination only needs an exporter. A new split-deployment transport needs an
exporter and a matching source. These additions would give users more choices
for connecting the monitoring pipeline to their existing systems while reusing
collectors, converters, and verdict services.

A Web UI is an important usability direction. It could discover the ROS graph,
let users select topic, service event, and action feedback/status streams,
message fields, transformers, and checking components, and show the placement
of runtimes across hosts. The UI could then validate the selected links, create
the JSON deployment description, invoke the configuration generator, and
present data records and verdicts.
This work would turn several manual YAML and deployment decisions into a guided
workflow. Its evaluation should study whether users can build correct monitoring
setups with less effort and whether the generated configuration still makes all
deployment and rule choices visible.

Finally, evaluation on physical robots and longer field runs would strengthen
the practical evidence. Such studies could compare integrated and split
placement under constrained robot hardware, measure the effect of high-rate
sensor data, introduce broker and network failures, and examine how operators
use the recorded evidence. These experiments would connect the reusable
infrastructure developed in this thesis with the wider field-based testing and
operational monitoring process.
